\section{Outline}
\label{sec:ch01:outline}

Section~\ref{ch01:sec:related_work} reviewed the literature along the main
conceptual gaps that motivate this monograph. Those gaps concern, in
particular, the mismatch between expressive temporal-task semantics and
computationally manageable multi-robot coordination, the difficulty of
maintaining task correctness under incomplete models and evolving requests, and
the further complications introduced by communication limits, uncertainty, and
dynamic task structure. The purpose of the present section is to turn that
literature map into the organizing logic of the book. Rather than revisiting
the earlier review, the discussion below explains how the later chapters take
up these issues in a deliberate order, starting from the nominal synthesis
setting and then moving toward progressively richer coordination scenarios.
Where appropriate, the chapter descriptions also identify the published papers
on which the corresponding technical developments are based.
The resulting structure is cumulative, but it is also driven by the gaps
identified in the previous section. Chapter~\ref{chap:preliminaries}
establishes the formal basis in the single-robot case, which provides the
technical reference point for the rest of the monograph. Chapter~\ref{chap:bottom-up}
addresses the difficulties that arise when coordination is induced by local
temporal tasks, partial knowledge, and execution-time collaboration.
Chapter~\ref{chap:top-down} considers the complementary setting in which the
mission is specified at the team level and must be decomposed and assigned
without losing its global temporal structure. Chapter~\ref{chap:communication}
then turns to the informational side of the same problem by studying limited
communication and online human involvement. Chapter~\ref{chap:uncertainty}
extends the framework further to probabilistic models, dynamic constraints, and
richer task semantics over evolving signals. Finally,
Chapter~\ref{chap:outlook} returns to the broader methodological picture and
highlights the open problems that remain at the intersection of formal
synthesis, scalable coordination, and adaptive autonomy.



\subsection{Single-robot foundations}
\label{sec:ch01:outline:foundations}

Chapter~\ref{chap:preliminaries} establishes the formal foundation on which the
rest of the monograph depends. The starting point is a finite abstraction of
robot motion over a partitioned workspace. This abstraction makes it possible
to describe robot behavior in symbolic terms while preserving the connection to
continuous dynamics through a lower-level control layer. Linear temporal logic
is then introduced as the task language used to express sequencing, safety,
recurrence, and related behavioral constraints. This abstraction--execution
view follows the single-robot motion-and-action planning framework in
\cite{guo2013motion} and the later human-in-the-loop ROS implementation in
\cite{baran2021ros}.
From there, the chapter develops the standard synthesis pipeline. A temporal
formula is translated into a B\"uchi automaton, combined with the motion
abstraction through a product construction, and searched for an accepting run
that yields a discrete plan. That plan is then interpreted by a hybrid
execution strategy that connects symbolic decisions to continuous motion.
Chapter~\ref{chap:preliminaries} therefore plays a dual role. It supplies the
technical preliminaries used throughout the book, but it also fixes the
nominal benchmark against which the later extensions should be read.

\subsection{Bottom-up coordination under local tasks}
\label{sec:ch01:outline:bottomup}

With that baseline in place, Chapter~\ref{chap:bottom-up} turns to multi-robot
coordination in a setting where the tasks are assigned locally rather than
globally. The emphasis is on how coordination emerges when robots that are
initially planned in a decentralized way begin to affect one another through
shared resources, action dependencies, missing knowledge, or temporary needs
for assistance. This perspective is especially natural for heterogeneous
teams, where the task of each robot is shaped by its own capabilities and by
the information available in its local neighborhood.
The first part of the chapter studies task relaxation and online adaptation in
partially known workspaces. The problem there is no longer only how to compute
a satisfying plan, but how to preserve meaningful task progress when the
workspace model changes during execution. This development follows the line of
work on plan revision, infeasible local specifications, multi-agent plan
reconfiguration, and knowledge transfer in
\cite{guo2013revising,guo2012infeasible,guo2015multi,guo2014distributed}.
The chapter then moves to dependent local tasks with collaborative actions,
where one robot can complete part of its task only if another robot provides
the required support. It closes with contingent collaboration, in which
requests are triggered during execution rather than specified exhaustively in
advance. These mechanisms are based on the bottom-up coordination,
heterogeneous task-and-motion coordination, and contingent formation-control
frameworks in~\cite{guo2015bottom,guo2016task,guo2017hybrid}. What
distinguishes the chapter as a whole is that coordination is not imposed by a
single team-level mission. It arises from local plans whose interactions become
visible only as execution unfolds.

\subsection{Top-down synthesis for global tasks}
\label{sec:ch01:outline:topdown}

A different organizational principle appears in
Chapter~\ref{chap:top-down}. Here the point of departure is not a collection of
local specifications, but a global temporal mission posed at the level of the
team. The central question is how such a mission should be decomposed and
distributed so that the resulting robot behaviors remain consistent with the
original team objective. In this sense, the chapter complements
Chapter~\ref{chap:bottom-up}. The earlier chapter studies how coordination can
emerge from local dependence, whereas the present one studies how it can be
derived systematically from global mission structure.
The chapter begins with partial-order-based task decomposition and assignment,
where the logical structure of the mission is separated from the choice of
which robot or subteam should implement each part. This part is tied to the
poset-product and online synchronization methods in
\cite{liu2022arxiv,liu2024time}. It then considers hierarchical coordination
for large-scale fleets, where abstraction across multiple planning layers is
used to contain computational growth without discarding mission-level
consistency. The final part addresses continual tasks, for which new demands
may appear while the team is already executing a plan; this development follows
the hierarchical continual-task framework in~\cite{luo2025hulk,wang2026hector}.
That setting brings the top-down perspective closer to online coordination, but the
organizing idea remains global from the outset. The mission is given at the
team level, and the synthesis problem is to allocate and refine it without
losing its temporal structure.

\subsection{Partial knowledge and limited communication}
\label{sec:ch01:outline:partial}

The next step is motivated less by task structure than by operational reality.
Even if the temporal objective is well posed and the coordination architecture
is sound, robot teams rarely act with complete and perfectly shared knowledge.
Chapter~\ref{chap:communication} addresses this difficulty by treating
information flow as part of the planning problem rather than as a background
assumption.
The chapter first examines settings in which communication and collaboration
must be planned jointly, following the adaptive communication and coordination
framework in~\cite{zhang2026cocoplan}. It then considers exploration and
coordination under intermittent communication, where the robots cannot rely on
continuous connectivity and must instead operate with delayed or sparse
information exchange. This part is connected to the earlier data-gathering,
temporal-task, team-wise intermittent-communication, and simultaneous
exploration and inspection developments in
\cite{guo2018multirobot,kantaros2019temporal,wang2026decentralized,chen2026slei3d}.
The last part studies online human interaction and request fulfillment, where
an operator may insert requests, guide decisions, or supply missing intent
without replacing the underlying autonomous planning framework, following the
human-oriented exploration and assistance frameworks in
\cite{jhunjhunwala2021ihero,zhang2025flykites,tian2026bodyguards}.
The common thread across these sections
is that temporal-task coordination must now coexist with incomplete awareness
of the environment, of peer plans, and of the future availability of
communication itself.

\subsection{Dynamic uncertainty and constraints}
\label{sec:ch01:outline:uncertainty}

Chapter~\ref{chap:uncertainty} pushes the framework further by weakening the
deterministic assumptions that still remain in the earlier chapters. In some
settings the difficulty lies in stochastic outcomes, in others in moving
constraints, evolving signals, or task semantics that cannot be expressed
adequately over static propositional labels alone. The chapter brings these
issues together under a broader view of temporal-task planning in dynamic and
uncertain environments.
The chapter opens with probabilistic temporal tasks and safe-return
constraints, where correctness must be interpreted alongside risk and
recoverability. This part follows the probabilistic temporal-task planning and
hierarchical safe-return formulations in
\cite{guo2018probabilistic,guo2023hierarchical}. It then studies reactive
coordination under dynamic temporal logic tasks, where the environment
is dynamic and can invalidate purely offline decisions
even when the original specification is sound, partially
based on results from~\cite{zhao2026umbrella}.
The final section considers collaborative optimal transport under signal
temporal logic tasks, thereby moving beyond the propositional setting that
dominates the nominal synthesis pipeline and connecting the chapter to the
coupled LTL/STL multi-robot framework in~\cite{lindemann2019coupled,zhang2026clot}.
By this stage, the monograph has moved well beyond the question of whether a
symbolic task can be satisfied on a fixed abstraction. The problem is now how
formal task semantics can still guide planning when the system model, the
environment, and even the relevant signals evolve during execution.

\subsection{Outlook and open problems}
\label{sec:ch01:outline:outlook}

Chapter~\ref{chap:outlook} closes the monograph by drawing out the research
questions that remain unresolved after the technical development of the
preceding chapters. The aim is not to append a detached list of future topics,
but to identify the main fault lines that cut across the entire book. These
include the scaling behavior of compositional and reactive synthesis,
coordination for continual tasks in nonstationary or adversarial environments,
and the integration of learning-based methods with formal guarantees that
remain meaningful at the system level. The discussion is framed by the recent
survey perspective in~\cite{li2025survey}, while remaining focused on the
specific methodological trajectory developed in this monograph.
For that reason, the final chapter should be read as a continuation of the
earlier material rather than as a separate epilogue. The challenges discussed
there arise directly from the limitations of the methods developed in
Chapters~\ref{chap:preliminaries} through~\ref{chap:uncertainty}. The overall
arc of the monograph is therefore intentionally open-ended. It develops a
coherent set of foundations and coordination frameworks, but it also makes
clear that the broader problem of hybrid coordination under complex temporal
tasks remains an active and expanding area of research.
